\subsection{Fortune’s Fool by Numbers}

\begin{quotex}
But I think about the hebdomads for myself, and I keep these speculations in my memory rather than sharing them with any of those people whose impertinence and insolence permits nothing to be analyzed without joking and laughter. Accordingly, don't be opposed to obscurities stemming from brevity; since they are the faithful guardian of the secret they have this advantage: they speak only to those who are worthy. Therefore, as is the customary practice in mathematics and also in other disciplines, I have first put forward terms and rules, in accordance with which I shall work out all that follows.
\end{quotex}


Boethius likely had a connection to esoteric doctrine via the Pythagoreans (his mundane connection to neo-Pythagoreanism is considered by scholars to be settled). He was a ``classically educated'' man, but as Miranda Lundy points out the four sciences of Number, Music, Cosmology \& Geometry were actually formerly \emph{sacred} sciences, and most definitively not a preparatory or prerequisite curriculum for anything else except sacral or transformative theology or philosophy. Sarah Lassin argues that Boethius' reference to the Egyptian ``hebdomads''\footnote{\url{http://individual.utoronto.ca/pking/translations/BOETHIUS.On\_the\_Hebdomads.pdf}} (which plays an important role within the Neopythagorean literature of Hierocles and Nicomachus of Gerasa\footnote{\url{http://books.google.com/books?id=dCELAAAAIAAJ\&pg=PA94\&lpg=PA94\&dq=nicomachus+hebdomads\&source=bl\&ots=TnzwHveyQF\&sig=ENK2-4LKLaMFDcxMhWkQdBwkqf0\&hl=en\&ei=wwPOTd6wN8PngQevmLihDA\&sa=X\&oi=book_result\&ct=result\&resnum=3\&ved=0CCoQ6AEwAg\#v=onepage\&q\&f=false}}) is highly significant for the discussion of God \& creation. In fact, it is not necessary even to mention the Egyptian origins of these terms, since Proclus discusses them\footnote{\url{http://books.google.com/books?id=z-M-AAAAYAAJ&pg=PA316&lpg=PA316&dq=hebdomads+triads&source=bl&ots=7lW956x_Vy&sig=77Nd93YLZNUyh5LvWP3foh-9FME&hl=en&ei=TKXMTZfrFYK4tgfn_u2WBQ&sa=X&oi=book_result&ct=result&resnum=1&ved=0CBYQ6AEwAA\#v=onepage&q=hebdomads\%20triads&f=false}} extensively on Plato, and Plato (and also Herodotus) thought that the Greeks derived much of their mysteries from Egypt. The Egyptian context\footnote{\url{http://www.sofiatopia.org/maat/ten_keys.htm}} makes the connections between philosophy and religion even more apparent:

\begin{quotex}
Hermes tells Tat (XIII), that ``the tent'' or ``tabernacle'' of the Earthly body was formed by the circle of the Zodiac (XIII.12 \& Ascl.35) and dominated by fate, who's decrees, according to the astrologers, were unbreakable. The seven planets represented the ``perfect movements'' of the Deities, the unalterable ``will of the Gods'' as expressed in predictable astral phenomena. Magicians tried to compel this will, while Hermetism did not try to resist fate, but irreversibly. The existence of the Deities was acknowledged (they belonged to the order of creation and were the object of sacrifices and processions and the celestial Powers ruling the astrological septet). Indeed, the Deities, Hermes and God were situated in the eighth, ninth and tenth sphere (Ogdoad, Ennead and Decad). The ``eighth'' involved purification, Self-knowledge and the direct ``gnostic'' experience of the ``Nous'' as ``logos'', whereas in the ``ninth'' man was deified by assuming God's attributes, as did the Godman Hermes, in particular His Universal Mind, the Divine Nous, Intellect or ``soul of God'' (XII.9). The ``tenth'' or Decad was God Himself for Himself.

\end{quotex}


And the Egyptians mysteries are linked to Hermes Trismegestus, which context resurfaces during the Renaissance when Greek translators made the Emerald Tablet available in Italy. This Egyptian link might also explain the persistence at Alexandria of a Christian esoteric tradition- Alexandria is also the home of the tragic conflict concerning Hypatia and the library, a conflict that did not occur until the Greek influence had waned in the Church.

Given that Boethius was the schoolmaster of medieval and early modern Europe, one might almost say that these doctrines ``sub-created'' Europe. Thomas Aquinas discusses\footnote{\url{http://individual.utoronto.ca/pking/translations/AQUINAS.Exposition\_of\_Hebdomads.pdf}} Boethius and the hebdomads in a treatise, \& it is instructive to compare the two discussions\footnote{\url{http://individual.utoronto.ca/pking/translations/BOETHIUS.On\_the\_Hebdomads.pdf}}; by Aquinas' time, the debate is far more abstract and limited to logical distinctions, as opposed to being ``transformative''.

As perennialist studies on Boethius\footnote{\url{http://www.studiesincomparativereligion.com/Public/articles/Boethius\_Three\_Musicians-by\_Joscelyn\_Godwin.aspx}} make clear, his theories of music and number (as found in minor tractates) gives a better background for the dialectic that occurs within theConsolation of Philosophy. Boethius was not simply a well-informed ``liberal'' who opposed barbarian elements at Court --- he was a lover of wisdom who was in an intimate and self-converting dialogue with Sophia herself, one who was willing to pay the final price for that which he had found and was finding\footnote{\url{http://www.hermes-press.com/philosophy\_emboldening0.htm}}.

\begin{quotex}
During the Dark Ages in Europe, acquaintance with the works of Plato was at second or even third hand, through the writings of authors such as Macrobius, Martianus Capella, Augustine, Boethius, Calcidius' translation of the \emph{Timaeus}, and John Scotus Erigena. The most important contribution in the vernacular was provided by the late ninth century reworking of Boethius' \emph{De consolatione Philosophiae}.
\end{quotex}


\emph{The Emboldenment of Philosophy} became one of the most popular books throughout the Middle Ages. It was translated into Old English by Alfred the Great, into Middle English by Chaucer, and into Elizabethan English by Queen Elizabeth.

As we see from Proclus' statement, the Neo-Platonists emphasized dialectical interchange in their transmutational teaching procedures within the Perennial Tradition--one of the important components of the Wisdom Teachings. However, the first Perennialist teacher who initiated and developed all facets of the Wisdom Teachings was Anicius Manlius Severinus Boethius (480-524 C.E.), including in his teachings and life experience each theme and emphasis we've outlined above.

While in prison awaiting a ghastly death, Boethius had experienced a definite inner dialectical interchange between his soul and Lady Philosophy (the spirit of the love of Wisdom). The narrative account in his \emph{Emboldenment of Philosophy}\footnote{\url{http://www.hermes-press.com/philosophy\_emboldening0.htm}} discloses how Boethius achieved transformation and self-understanding while communing with Lady Sophia in an inner spiritual domain. Boethius' book combined verse and alternating dialogue between himself and Lady Philosophy, organizing his narrative into different stages of his spiritual healing and transmutation. Lady Philosophy communes with Boethius in his inner being, bringing into the dialectical interchange others who were true devotees of philosophy: those who love and seek Wisdom.

The wisdom school of the Pythagoreans was influential enough that Augustine made an appeal to those who ``knew'' to the effect that the Hebdomads symbolized the perfection which was to become the Church, which was to rely in almost exclusive fashion\footnote{\url{http://orthodoxnotes.com/wp-content/uploads/downloads/2010/06/Liturgical-Practice-in-the-Fathers.pdf}} upon the mystery religions, perhaps even more than on Judaism, for the new doctrines. Christianity began as a conversion of existing elements, rather than a superimposition. Augustine's appeal is a sincere one.

Given that Boethius' interests in music and number coincide with Platonist streams of thought in Proclus and Plotinus in the Pythagorean mysteries, it is necessary to think that the esoteric link which created the Middle Ages through Boethius could be located in the harmonics and sacred geometry\footnote{\url{http://www.sacred-texts.com/eso/sta/sta19.htm}} of the celestial ``music of the spheres''. In order to understand the Middle Ages, we should ``close thy scholastics, and open thy Timaeus''. Boethius' background and writings, embodied in the \emph{Consolation}, are an excellent place to re-start.



\flrightit{Posted on 2011-05-14 by Logres}

\begin{center}* * *\end{center}

\begin{footnotesize}
\begin{sffamily}
\texttt{logres on 2011-05-14 at 22:39 said:}

This is a good example of impasse between ``orthodox'' and ``gnostics'':

\url{http://www.interfaith.org/forum/livergoods-esoterism-10888.html}

Trajan's orders seem reasonable from a political standpoint -- the Christian response was muted and in accord with Romans 13. However, it is clear that persecutions of the faith quickly became excessive. This probably gutted much of the leadership. If the Corpus Hermeticum dates from the 3rd century, then it is the obverse of the same coin which Clement was minting.

The chance was missed. And is still being missed. That is why Boethius is so important. In spiritual war, as in physical, it is often the best that are lost. To tie into Cologero's posts, there is still an opportunity perhaps even today to refashion Romanity as a vehicle, in the Catholic Church. If the chance is missed again, we will simply behold a repetition of the pattern, until what is meant to happen, what is already one, is manifest. From the debate linked above, it's clear that Mouravieff's Gnosis is an important tract.

\hfill

\texttt{Spock on 2011-05-15 at 09:40 said:}

logres, you say there is still a chance to fix the church, well yes, we can say this forever, but the truth is that leading perennialists have had continuing dialogues with the recent popes and it changes absolutely nothing because the church is and will only be a political vehicle which panders to its supporters... none of whom are even remotely interested in gnosis and tradition. What they love are superficial rituals and the joy of submission to authority, because they do not love wisdom, they would rather be told and given things like children rather than do the task themselves. That in itself in contrary to the active gnostic life. If one really wanted to change the church he must first change the people, but that is the impossible task in this the last phase of the kali yuga.

\hfill

\texttt{James O'Meara on 2011-05-15 at 10:10 said:}

Freke and Gandy quote modern scholarship to show that the persecutions were quite minor, both in time 5 years, total and numbers 10 killed in Alexandria! ; Origen is quoted to the effect that the number of `martyrs' was `easily counted'.

\url{http://tinyurl.com/5v7w5nv}


The ``great persecutions'' could hardly have `gutted much of the leadership'' although the orthodox literalists did a pretty good job decimating the Gnostics and other `heretics'.

As always, the primitive Christians seem to specialize in self-serving propaganda and literary fakery the letters of Jesus to Pilate, anyone? . Speaking of which, interestingly, the Gnostics rejected the ``pastoral epistles'' as spurious, just as modern critics do, while Irenaeus insisted on them, showing his unerring eye for self-serving fakery, and proving that Guenon was wrong to dismiss Biblical criticism as a modernist evil. Even today, listening to radio preachers, it's amazing how much of their bible-thumping comes from Timothy, Titus and Peter, plus a whole heaping helping of Apocalypse. ``Modernist criticism'' would disarm a lot of the Protestantism Guenon opposed.

\hfill

\texttt{logres on 2011-05-16 at 08:58 said:}

It wouldn't surprise me to find out the persecutions were exaggerated -- however, I can't find footnote 138 (and those next to it), to see which scholars they are citing. The argument doesn't really depend upon who killed whom, but rather upon ``elective affinities''. Certainly, it is possible that the Church became exactly what She portrayed Rome as being. Tomberg, Mouravieff, etc. aren't trying to esoterically ``capture'' the exoteric Church as it has developed (which is splitting up all on its own into fragments) but to reunite those elements which should be (and are) One. So I don't think it matters what the overly visible Church does. At least not for purposes of working in the areas of Form.

However, it certainly matters to set the record straight, if what one means by that is to revisit assumptions which one thinks one knows, and for the purposes of refuting aggressive error. Thanks for drawing attention to this issue.

\hfill

\texttt{logres on 2011-05-16 at 09:04 said:}

The Church (in other words) is not only the people, and perhaps in a sense not the people at all. A difficult task, I know.

\hfill

\texttt{logres on 2011-05-18 at 22:03 said:}

B16 speech on Proclus \& Pseudo-Dionysius:\\
\url{http://www.zenit.org/article-22588?l=english}


\hfill

\texttt{Dimitri on 2022-06-30 at 19:18 said:}

Though it has been over a decade since his comment was posted, I would still like to correct Mr. O'Meara's many errors (which stem in large part from his relying on the absurdities of the so-called ``Christ myth theorists'' as a source). In his first point, he claimed the persecutions of early Christians were exaggerated. This is somewhat true to an extent, though they were nowhere near ``minor'' like he claimed. He gives the impression that there were perhaps a few dozen people killed and that the overall period of persecution was very short, but if only he had done his research he would have seen just how wrong he is. How he arrived at the idea that persecution only lasted for ``5 years total'' is beyond me, since even if we were to only include imperially ordered persecution, we would still arrive at a total of 13 years of persecution under Decius, Valerian, Diocletian and Galerius. And this is only if we exclude the many local persecutions where Christians were attacked by regular people without any official orders, or were persecuted under various governors whose authority was limited to specific regions of the empire. For the first century or so after the beginning of Christianity, Christians were fairly unpopular with the people and had to avoid the public sphere. But beyond this, O'Meara also greatly downplays the death toll even from the imperial persecutions. The Diocletianic Persecution alone claimed well over 3000 lives, even according to ``modern scholarship'' (and keep in mind that the number of Christians around then was smaller than it is today). There were many thousands of Christians martyred for their religion throughout the years, his deceptive use of Origen's statements notwithstanding.

As for whether ``\,`modernist criticism' would disarm a lot of the Protestantism Guenon opposed'', perhaps this is true, but it certainly doesn't replace it with anything worthwhile or meaningful. In fact, it seems to have made the situation worse, since we're now at a place where regular people (at least in the West, and this is especially true of Americans) see Christianity only through the lense of fundamentalism and mistakenly think that to be Christian they MUST believe in a strictly literal interpretation of all Scripture (among other errors) and have therefore forgotten the essence of the religion entirely. Perhaps the reason Guenon opposed it was because these ``critics'' and ``scholars'' do not really understand the contents or the implications of what they are supposed to be studying? For example, Cologero recently made a post about a scholar discussing how the Gospels were authored by a Christian elite (which would of course go against the ``traditional attribution'' to ``early eye-witnesses''), and yet from my experience this information would sooner be used in a polemical attack (as O'Meara himself tried to do here) by those who in their ignorance possess a vigorous contempt for what they think is Christianity than it would be to actually learn anything useful.

\hfill

\texttt{Dimitri on 2022-06-30 at 21:02 said:}

Actually, I realize now I forgot to mention the fact that ``modernist criticism'' itself only exists because of Protestantism and its nontraditional elements, such as its tendency for secularization and adherence to various profane ``philosophies'', and so it could hardly be an antidote to it. So not only does it not replace it with anything meaningful as I said before, but ironically it's similar to an Ouroboros, in that it's like a snake consuming itself.

\end{sffamily}
\end{footnotesize}

